% 2.6 Autonomous Edge-Cloud Orchestration (State of the Art)
\subsection{Autonomous Edge-Cloud Orchestration}
\label{sec:ch2_sota_agent}

Section~\ref{sec:ch2_intelligent_orchestration} established the theoretical foundations for \ac{AI}-driven orchestration, including the three-stage orchestration problem, the \ac{PARES} capability framework, and the mathematical benefits of multi-agent coordination. This section evaluates how existing systems implement these concepts, reviewing: (i) \ac{LLM}-based orchestration approaches for intent processing and code generation, (ii) agent-based orchestration systems spanning domain-specific applications, multi-agent coordination, and infrastructure management, and (iii) fundamental limitations that motivate \ac{AgentEdge}.

\subsubsection{LLM-based Orchestration}

\ac{LLM}-based approaches divide into two primary categories: intent processing and code generation.

\textit{Intent Processing and Translation:} Translating natural language intents into network configurations requires breaking down ambiguous requests into verifiable actions. Recent work~\cite{intent_driven_ran2021} uses few-shot learning to break down and validate intent in next-generation networks. Other studies~\cite{IntentBasedNetworkManagementDRCN2024} show intent extraction capabilities for \ac{6G} core networks. Advanced approaches use a four-stage methodology to convert natural language intentions into Network Service Descriptors using Code Llama \ac{LLM}~\cite{LLMIntentServiceConfig2024}, incorporating few-shot learning, similarity-based retrieval, closed-loop validation, and human feedback for continuous improvement.

\textit{Code Generation and Configuration Management:} Beyond intent extraction, automating network configuration requires generating correct code while addressing privacy and error-handling concerns. Recent studies~\cite{EnhancingNetworkManagementMSR2023} present task-specific code generation for network graph analysis and manipulation, combining application-specific prompts with program synthesis techniques to address explainability, scalability, and privacy challenges. Complementary work~\cite{IntentNetworksLLMFramework2024} introduces a modular framework that addresses data privacy and \ac{LLM} errors in router and firewall configurations.

However, current \ac{LLM}-based orchestration systems face critical limitations~\cite{NEURIPS2024_6aebba00, struggle_with_real_time_adaptation2025}: model hallucinations that generate incorrect configurations, inference latencies that prevent real-time adaptation, insufficient tool integration for infrastructure \acp{API}, and monolithic architectures unsuitable for distributed edge-cloud environments. The monolithic nature of these systems creates a fundamental mismatch with the distributed requirements of edge-cloud orchestration.

\subsubsection{Agent-based Orchestration}

Building upon these \ac{LLM} foundations, agent-based orchestration systems represent the next evolution, using autonomous \ac{AI} agents for coordinated service orchestration across heterogeneous environments.

\textit{Domain-Specific Agent Approaches:} Implementations target specialized networking applications with customized agent architectures. In wireless network management, researchers~\cite{WirelessAgentLLM2025} demonstrate intelligent network management through \ac{LLM}-powered agents. For fault management, other work~\cite{LLMAgentsProactiveFaultTolerant2024} introduces an \ac{LLM}-based agent system with integration of memory, knowledge base, and network infrastructure. This system autonomously manages the complete fault management cycle through proactive traffic monitoring, delay detection, and flow rerouting. For network planning applications, researchers~\cite{RadioMapAgentsNetworkPlanning2025} present a three-module framework with specialized tools for geographic map retrieval, environment creation, and automated radio map generation. However, these approaches rely on single-agent architectures that cannot distribute reasoning across specialized roles, lack formal capability definitions, and execute actions directly without pre-validation in simulated environments.

\textit{Agent Coordination and Collaboration:} Moving beyond single-agent approaches, researchers have developed frameworks for multi-agent cooperation and coordinated decision-making. Recent work~\cite{HermesLLMFramework2024} uses a chain-of-agent approach that separates network modeling into specialized roles using \ac{YAML}-based blueprints with mathematical formulas and pseudocode. Other approaches~\cite{MaestroLLMCollaborative2024} present multi-agent systems with \ac{LLM}-based agents for conflict resolution and resource negotiation on the business, service, and network planes. Additional work~\cite{xiao2025agentnet} uses agents across the application, physical, and network layers, evaluated through industrial automation and virtual reality applications. These frameworks typically use flat coordination structures without hierarchical organization, accept \ac{LLM} outputs without constraint verification, and test plans through trial-and-error on production infrastructure rather than simulation.

\textit{Infrastructure Orchestration and Service Management:} Finally, infrastructure orchestration approaches focus on agentic implementations for distributed service management and deployment. Recent work~\cite{DataConnectorAgentFramework2024} uses machine learning models for \ac{QoS} predictions and allows autonomous migration of containerized applications in Kubernetes-based edge-cloud environments. Other work~\cite{service_orch_2025} presents an intent-based orchestration approach using four specialized agents with distributed coordination. This system employs Model-Context Protocol integration for tool standardization and uses fine-tuned Mistral-7B and Qwen3-0.6B models through \ac{QLoRA}~\cite{dettmers2023qlora} for Kubernetes configuration tasks. These systems assume cloud-centric architectures with reliable high-bandwidth connectivity, limiting applicability to heterogeneous edge-cloud environments, and address deployment but not migration, scaling, or power management. Despite their varied approaches, these agent-based systems share underlying assumptions that constrain their applicability to distributed edge-cloud scenarios.

\subsubsection{Fundamental Limitations and Research Gaps}

The \ac{LLM}-based and agent-based approaches reviewed above share critical gaps for distributed edge-cloud environments. They focus on centralized cloud infrastructures, target narrow applications without addressing complete service lifecycle management, and lack formal agent capability definitions and coordination mechanisms for distributed multi-agent collaboration.

Beyond domain-specific systems, generic agentic approaches from the broader \ac{ML} literature~\cite{yao2023react, LATS} such as \ac{LATS} employ trial-and-error learning requiring direct experimentation on production infrastructure. This creates reliability concerns for deployments where failures cascade across distributed services. In contrast, traditional optimization methods offer deterministic guarantees but have the converse limitation: they solve placement problems given formally specified objectives, yet cannot interpret natural language intents or handle partial infrastructure states~\cite{6g_struggles_orchestration2025}. These gaps reveal a persistent safety-autonomy trade-off: autonomous agents offer adaptability but risk production failures, while rule-based systems provide safety but lack flexibility.

\subsubsection{Addressing the Limitations with AgentEdge}

\ac{AgentEdge} addresses this safety-autonomy trade-off through its \ac{ActSimCrit} methodology, where agents Act in simulation, observe outcomes through Simulation, and receive feedback through Critique before production deployment. Instead of choosing between unsafe exploration (\ac{RL}/agents) or rigid heuristics (rule-based), \ac{AgentEdge} introduces a third path: validated autonomy. By integrating a high-fidelity digital twin directly into the agent's reasoning loop, it enables risk-free exploration of complex orchestration strategies. This allows agents to learn from failure in simulation rather than production, combining the adaptability of generative AI with improved reliability through pre-deployment validation. Chapter~\ref{ch:distributed-agentic-ai} presents this framework, demonstrating how distributed specialized agents can safely orchestrate mission-critical \ac{6G} infrastructure.
